\documentclass[10pt,a4paper]{amsart}
\usepackage[margin=19mm]{geometry}
\usepackage{amsmath,amssymb,mathtools}
\usepackage[scr=rsfs,cal=euler]{mathalfa}
\usepackage{xfrac}
\usepackage[unicode,hidelinks,bookmarks=false]{hyperref}
\newcommand{\ie}{i.e.\ }
\newcommand{\cf}{cf.\ }
\newcommand{\resp}{respectively\ }
\newcommand{\Subsection}{\subsection}
\providecommand{\nobreakdash}{\nobreak\hbox{-}}
\setlength{\emergencystretch}{2em}
\multlinegap0pt
\usepackage{xcolor}
\usepackage{bm}
\usepackage{caption}
\usepackage{amssymb}
\usepackage{mathtools}
\usepackage{tikz}
\usepackage{tensor}
\usepackage[shortlabels]{enumitem}
\newtheorem{prop}{Proposition}
\def\Gt{\widetilde{G}}
\def\Mt{\widetilde{M}}
\def\Rt{\widetilde{R}}
\def\xt{\tilde{x}}
\def\yt{\tilde{y}}
\title{A Hamiltonian approach for point vortices on non-orientable surfaces}
\author{Nataliya A. Balabanova, James Montaldi}
\date{Source: arXiv:2202.06160v3 (CC BY 4.0)}
\begin{document}
\maketitle

\noindent\emph{Source and licence.} A Hamiltonian approach for point vortices on non-orientable surfaces. Version arXiv:2202.06160v3 (CC BY 4.0), distributed by arXiv under the Creative Commons Attribution 4.0 International licence, http://creativecommons.org/licenses/by/4.0/, as recorded in the arXiv licence metadata for this exact version and shown on the abstract page https://arxiv.org/abs/2202.06160v3. Full licence notice: \texttt{licenses/arxiv-2202.06160-CC-BY-4.0.txt}.\par\smallskip

\noindent\textit{Editorial scope.} Counted unit: the article's prop function (source0.tex lines 666--669). The article's complete original proof of it is delivered with it (source0.tex lines 670--704, one \texttt{proof} environment of 35 lines). No mathematics is written, repaired or re-derived: the statement and the proof are the article's own lines, in the article's own notation, and no retained line is substituted or re-flowed.  Retained uncounted as the source-local context the proof needs: lines 175--178; lines 657--660.  Nothing was omitted from any retained span, so the package carries no omission record.  No retained line contains a citation, so the wrapper carries no bibliography and no bibliography entry.  No retained line contains a figure environment, an image-inclusion command or drawing code, so no image is delivered and the package declares no asset dependency.  Editorial formatting only, in addition: an AMS \texttt{amsart} preamble that restates the article's own declarations and macros used by the retained lines, namely the environments \texttt{prop} on separate counters, together with the standard packages those lines need.  The retained environments print in this document as Proposition 1 in order of appearance, whereas the article prints the same items with the numbering of its own full document.  The complete native source of the article as distributed in the arXiv e-print for this exact version is bundled unchanged as \texttt{sources/122012/source0.tex}.
\begin{equation}
\label{eq: Robin function definition}
R(x) = \lim_{y\to x} \left(G(x,y) -\frac1{2\pi}\log d(x,y)\right).
\end{equation}

\begin{equation}
\label{eq: twisted Robin}
     R_M(x) = \Rt(\xt)-\Gt(\xt,\tau(\xt))
\end{equation}

\begin{prop}
\label{st: green function and robin function}
 $G_M$ is indeed  a twisted Green's function of the Laplacian on $M$, and $R_M(x)$ is the Robin function of $G_M(x,y)$.
\end{prop}

\begin{proof}


Consider the Laplacian of $G_M$ as a function on $\Mt$:
\[
\Delta_{\xt}\left( \Gt(\xt,\yt) -\Gt(\xt,\tau(\yt))\right) = \delta_{\yt}(\xt) - \delta_{\tau(\yt)}(\xt)  
\]
For a fixed $y$ (and, consequently, $\yt$) $\Delta_{\xt}\left(\Gt(\xt,\yt) -\Gt(\xt,\tau(\yt))\right)$ descends to $M$ as a twisted function $\widetilde{\delta}_y(x)$. \textcolor{black}{Note that $\widetilde{\delta}_y(x)$ is a \textit{twisted function} defined on the manifold $M$ (whereas $\delta_{\yt}(\xt) - \delta_{\tau(\yt)}(\xt)$ is a function on $\Mt$).} 

Analogously to the regular Dirac delta function, it indicates when one of the preimages of $x$ on $\Mt$ coincides with one of the preimages of $y$ on $\Mt$.  

As one can easily see, $G_M$ depends on the choice of the preferred lift, since it is twisted in both of its variables. 

Now we want to check that $R_M$ coincides with the desingularization of $G_M$.

Saying  that $x\to y$ for two points $x,y\in M$ is equivalent to one of the two statements for $\Mt$: $\xt\to \yt$ or $\xt\to \tau(\yt)$. Observe, however, that  $R_M$, despite being a limit, is a function of one variable only and is therefore covered by a function on $\Mt$. 

 Thus, we are restricted to one lift only: either $\xt\to \yt$ or $\xt\to \tau(\yt)$, depending on $\tau$ and our choice of preferred lifts. Since $\xt$ is a `mute' variable, it does not matter which of the two options we choose, as can be seen from the calculations  below.






We have from (\ref{eq: Robin function definition}) and (\ref{eq: twisted Robin}):
\begin{equation}
\label{eq: Robin form HAM}
    \begin{split}
        \Rt(\yt)- \Gt(\yt,\tau(\yt)) 
        =& \lim_{\xt\to \yt}\left(\Gt(\yt,\xt) -\frac{1}{2\pi}\log d(\yt,\xt)  - \frac12\Gt(\yt, \tau(\xt))  - \frac12\Gt(\xt,\tau(\yt))\right)  \\ =&\lim_{\xt\to \yt}\biggl(\frac12\Gt(\yt,\xt) + \frac12\Gt(\tau(\yt),\tau(\xt))-\frac{1}{2\pi}\log d(\yt,\xt)  \\& \qquad\quad -\frac12\Gt(\yt, \tau(\xt))  - \frac12\Gt(\xt,\tau(\yt))\biggr)  \\ =&\lim_{\xt\to \yt}\left(G_M(\yt,\xt) -\frac{1}{2\pi}\log d(\yt,\xt) \right).
    \end{split}
\end{equation}
Since we choose one lift out of two, the 
explicit form of $\Rt(\tau(\yt)) - \Gt(\yt,\tau(\yt))$  is the limit with $\tau(\xt)\to\tau(\yt)$, which can be shown to coincide with (\ref{eq: Robin form HAM}). Hence, $R_M(x)$ is a regular function on $M$, independent of our choice of preferred lift.
\end{proof}

\end{document}
